\chapter{Main Results}
\label{chapter:MainResults}

\section{Technical Development}

For reproducibility purposes, the code, configuration files, and generated
result artefacts used in this chapter are included in the project repository
submitted together with the thesis. At the time of writing, no public archival
link had been created for the repository, so reproducibility is supported
through the submitted codebase and the exported experiment outputs.

The implementation is organized around a typed experiment configuration, a
single Python entry point, and an export step that translates the trained
decision tree into deterministic MQL5 logic. This structure makes it possible
to reproduce the complete experiment from one configuration file while still
keeping the deployed strategy readable.

\subsection{Experiment Configuration}

The full experiment is controlled through one nested JSON configuration file.
The most important parts of the thesis configuration are the market-data
window, the indicator parameters, the tree settings, and the backtest
assumptions:

\begin{lstlisting}[language={},caption={Simplified thesis configuration for the XAUUSD H1 experiment.}]
{
  "results_dir": "results_xauusd_h1_3y",
  "experiment": {
    "market_data": {
      "symbol": "XAUUSD",
      "timeframe": "H1",
      "utc_from": "2022-01-01T00:00:00Z",
      "utc_to": "2024-12-31T23:00:00Z"
    },
    "tree": {
      "enabled": true,
      "max_depth_grid": [2, 3, 4, 5, 6],
      "min_samples_leaf_grid": [5, 10, 20, 50],
      "holdout_ratio": 0.2,
      "time_series_splits": 3,
      "confidence_bin_count": 5
    },
    "backtest": {
      "cost_bps": 5.0,
      "take_profit_pct": 0.06,
      "stop_loss_pct": 0.03
    }
  }
}
\end{lstlisting}

The full configuration file fixes the exact \texttt{XAUUSD} H1 experiment
discussed in this chapter, while the listing above highlights the most
important thesis parameters. Using one canonical configuration file reduces
ambiguity between the reported thesis setup and the executed code.

\subsection{Running the Experiment}

Once credentials are supplied through environment variables, the complete
experiment can be reproduced with one command:

\begin{lstlisting}[language=bash,caption={Running the thesis experiment from the command line.}]
python -m src.experiment \
  --config configs/run_config.xauusd_h1.json
\end{lstlisting}

This command retrieves the market data from a running MT5 terminal, constructs
the indicator features, generates candidate signals, trains the decision tree,
partitions the test signals into confidence bands, and writes the output
artefacts. The most important artefacts are \texttt{metrics.csv},
\texttt{trades.csv}, \texttt{rates.csv}, \texttt{equity\_curves.csv},
\texttt{confidence\_scores.csv}, and \texttt{config\_used.json}.

\subsection{Export to MQL5}

After training, the final decision tree is exported as deterministic
\texttt{if}-\texttt{then} logic in MQL5. In simplified form, the exported code
looks as follows:

\begin{lstlisting}[language=C++,caption={Simplified structure of the exported MQL5 decision tree.}]
double EvaluateDecisionTreeProbability(...)
{
   if (rsi <= 42.5)
   {
      return 0.1913043478;
   }
   else
   {
      if (buy_vote_count <= 3.5)
      {
         return 0.3919597990;
      }
      else
      {
         return 0.5161290323;
      }
   }
}
\end{lstlisting}

The exact thresholds vary with the trained model, but the important point is
that the exported tree remains transparent and directly inspectable. This is
consistent with the goal of the thesis: to use machine learning as an
interpretable signal-evaluation layer rather than as an opaque forecasting
model.

\section{Experimental Results}

The final experiment was conducted on \texttt{XAUUSD} H1 data requested from
1 January 2022 to 31 December 2024. Because the market is closed on weekends,
the retrieved dataset spans from 3 January 2022, 01:00 UTC, to 31 December
2024, 23:00 UTC and contains 17,749 hourly bars. The decision tree was trained
on development-period candidate signals and evaluated on a held-out
out-of-sample segment beginning on 28 May 2024, 14:00 UTC.

\begin{table}[htbp]
\centering
\caption{Experimental setup for the final XAUUSD H1 run.}
\label{tab:xauusd-h1-setup}
\small
\begin{tabular}{p{0.34\textwidth}p{0.52\textwidth}}
\toprule
Item & Value \\
\midrule
Instrument & \texttt{XAUUSD} \\
Timeframe & H1 \\
Requested period & 2022-01-01 to 2024-12-31 (UTC) \\
Retrieved observations & 17,749 hourly bars \\
Observed data span & 2022-01-03 01:00 to 2024-12-31 23:00 (UTC) \\
Development candidate signals & 1,148 \\
Out-of-sample candidate signals & 288 \\
Out-of-sample window & 2024-05-28 14:00 to 2024-12-31 23:00 (UTC) \\
Selected tree parameters & max depth = 3, min leaf size = 50 \\
Confidence grouping & 5 bands derived by quantile binning \\
\bottomrule
\end{tabular}
\end{table}

Table~\ref{tab:xauusd-h1-setup} summarizes the final experimental setup. The
important point is that the thesis results are based on one fixed configuration
and one fixed out-of-sample period. This reduces flexibility in the evaluation
stage and makes it easier to interpret the confidence-band results as a
comparison within a controlled setting.

\begin{table}[htbp]
\centering
\caption{Out-of-sample baseline performance on XAUUSD H1.}
\label{tab:xauusd-h1-baselines}
\footnotesize
\begin{tabular}{lrrrrrr}
\toprule
Strategy & Return & Sharpe & Max DD & Trades & Win rate & Profit factor \\
\midrule
EMA & 0.59\% & 1.65 & -0.27\% & 72  & 43.06\% & 1.51 \\
RSI & 0.79\% & 0.74 & -1.50\% & 210 & 38.57\% & 1.26 \\
MACD & 0.53\% & 0.86 & -1.01\% & 186 & 37.10\% & 1.18 \\
DMI & 1.64\% & 1.59 & -1.36\% & 282 & 40.43\% & 1.36 \\
KST & 0.16\% & 0.37 & -0.84\% & 152 & 34.87\% & 1.07 \\
MFI & 0.66\% & 0.50 & -2.08\% & 308 & 36.36\% & 1.14 \\
Multi-indicator & 1.27\% & 1.06 & -1.91\% & 241 & 40.25\% & 1.35 \\
Buy-and-hold & 11.80\% & 1.42 & -8.79\% & 0 & --- & --- \\
\bottomrule
\end{tabular}
\end{table}

Table~\ref{tab:xauusd-h1-baselines} shows that the unfiltered multi-indicator
strategy remained profitable out of sample, with a total return of 1.27\%, a
Sharpe ratio of 1.06, and a profit factor of 1.35. Among the single-indicator
baselines, the DMI strategy achieved the strongest out-of-sample total return,
while buy-and-hold unsurprisingly produced the largest absolute return over the
same market period. However, buy-and-hold also incurred a substantially larger
drawdown than the active trading strategies. This makes the multi-indicator
baseline a reasonable comparison point for the confidence-filtered variants.

\begin{table}[htbp]
\centering
\caption{Out-of-sample performance by decision-tree confidence range.}
\label{tab:xauusd-h1-confidence}
\footnotesize
\begin{tabular}{lrrrrrrr}
\toprule
Range & Min & Max & Return & Sharpe & Trades & Win rate & Profit factor \\
\midrule
Range 1 & 0.1800 & 0.1913 & 0.01\%  & 0.04  & 66 & 33.33\% & 1.00 \\
Range 2 & 0.2632 & 0.2632 & -0.44\% & -2.47 & 27 & 14.81\% & 0.35 \\
Range 3 & 0.3621 & 0.3864 & 0.14\%  & 0.59  & 27 & 40.74\% & 1.38 \\
Range 4 & 0.3920 & 0.3920 & 1.82\%  & 2.45  & 53 & 73.58\% & 5.57 \\
Range 5 & 0.4348 & 0.5161 & -0.24\% & -0.59 & 68 & 30.88\% & 0.89 \\
\bottomrule
\end{tabular}
\end{table}

Table~\ref{tab:xauusd-h1-confidence} contains the central result of the thesis.
The confidence bands do not form a perfectly monotonic ordering. In
particular, Range~5 did not outperform Range~4 even though it corresponds to a
higher predicted probability interval. Nevertheless, the table still shows that
the decision-tree confidence score was useful in this experiment for isolating a
stronger subset of candidate trades. Range~4 outperformed the unfiltered
multi-indicator baseline on total return, Sharpe ratio, win rate, and profit
factor, while Ranges~2 and~5 performed poorly. This indicates that the
confidence score was informative in the selected setting, but not uniformly
increasing in usefulness across all bins.

\begin{figure}[htbp]
\centering
\begin{tikzpicture}
\begin{axis}[
    width=0.88\textwidth,
    height=0.50\textwidth,
    xlabel={Out-of-sample H1 bar index},
    ylabel={Equity},
    xmin=1,
    grid=both,
    grid style={line width=.1pt, draw=gray!20},
    major grid style={line width=.2pt, draw=gray!35},
    legend style={
        at={(0.98,0.02)},
        anchor=south east,
        legend columns=1,
        draw=none,
        font=\footnotesize
    },
    tick label style={font=\small},
    label style={font=\small},
]
\addplot[
    very thick,
    color=blue!70!black,
] table[
    col sep=comma,
    x=bar_index,
    y=equity_multi_indicator_vote_oos,
] {figures/thesisFigures/data/xauusd_h1_3y_equity_oos.csv};
\addlegendentry{Multi-indicator vote}

\addplot[
    very thick,
    color=red!75!black,
] table[
    col sep=comma,
    x=bar_index,
    y=equity_tree_confidence_range_4_oos,
] {figures/thesisFigures/data/xauusd_h1_3y_equity_oos.csv};
\addlegendentry{Tree confidence range 4}

\addplot[
    thick,
    color=black!65,
    dashed,
] table[
    col sep=comma,
    x=bar_index,
    y=equity_buy_and_hold_oos,
] {figures/thesisFigures/data/xauusd_h1_3y_equity_oos.csv};
\addlegendentry{Buy and hold}
\end{axis}
\end{tikzpicture}

\caption[Out-of-sample equity curves]{Out-of-sample equity curves for buy-and-hold, the multi-indicator baseline, and Range 4.}
\label{fig:xauusd-h1-equity}
\end{figure}

Figure~\ref{fig:xauusd-h1-equity} compares the out-of-sample equity paths of
the main baseline, buy-and-hold, and the strongest decision-tree confidence
band. The buy-and-hold curve ends highest in absolute terms because the gold
market appreciated over the evaluation period, but it also experienced a larger
drawdown than the active trading strategies. The important comparison for the
research question is therefore between the unfiltered multi-indicator baseline
and Range~4. In this comparison, Range~4 ends higher than the unfiltered
baseline and is consistent with the higher Sharpe ratio, win rate, and profit
factor reported in Table~\ref{tab:xauusd-h1-confidence}.

\begin{figure}[htbp]
\centering
\begin{tikzpicture}
\begin{axis}[
    ybar,
    bar width=9pt,
    width=0.96\textwidth,
    height=0.53\textwidth,
    ymin=-1,
    ymax=80,
    ylabel={Percent},
    symbolic x coords={Range 1,Range 2,Range 3,Range 4,Range 5},
    xtick=data,
    xticklabel style={font=\small},
    yticklabel style={font=\small},
    label style={font=\small},
    legend style={
        at={(0.5,0.98)},
        anchor=north,
        legend columns=2,
        draw=none,
        font=\footnotesize
    },
    nodes near coords,
    every node near coord/.append style={font=\scriptsize, rotate=90, anchor=west},
    enlarge x limits=0.15,
    grid=major,
    grid style={draw=gray!20},
]
\addplot[
    fill=blue!55,
    draw=blue!80!black,
] table[
    col sep=comma,
    x=range_label,
    y=total_return_pct,
] {figures/thesisFigures/data/xauusd_h1_3y_confidence_ranges.csv};
\addlegendentry{Total return (\%)}

\addplot[
    fill=red!45,
    draw=red!80!black,
] table[
    col sep=comma,
    x=range_label,
    y=win_rate_pct,
] {figures/thesisFigures/data/xauusd_h1_3y_confidence_ranges.csv};
\addlegendentry{Win rate (\%)}
\end{axis}
\end{tikzpicture}

\caption[Confidence-range return and win rate]{Out-of-sample total return and win rate across the five confidence ranges.}
\label{fig:xauusd-h1-confidence-plot}
\end{figure}

Figure~\ref{fig:xauusd-h1-confidence-plot} makes the non-monotonic pattern more
visible. Range~4 stands out clearly, while Range~2 is the weakest group and
Range~5 falls back below both the baseline and the best confidence band. This
means that the decision-tree confidence should not be interpreted as a simple
``higher is always better'' score. Instead, in this experiment it functioned as
a useful but imperfect filtering signal that isolated one particularly strong
band of candidate trades from weaker ones.

\section{Summary of Main Findings}

Three findings are especially important. First, the unfiltered multi-indicator
strategy remained profitable in the held-out market segment, which means the
baseline itself was not degenerate. Second, the decision-tree confidence score
was useful for separating out one stronger out-of-sample band, because the
strongest band materially outperformed the baseline on several core metrics.
Third, the
relationship was not strictly monotonic across all five bands. The most
defensible conclusion is therefore not that higher confidence always implied
better trading performance, but that the confidence score contained useful
information for filtering candidate signals in the selected \texttt{XAUUSD} H1
setting.
